\chapter{マッチャーを定義しよう}\label{matcher}

本章は，ユーザーが自身でマッチャーを定義する方法を解説する．　

\section{\texttt{inductive pattern}宣言}\label{matcher-inductive}

マッチャーを定義する際には，そのマッチャーが処理するパターンコンストラクタをあらかじめ\emph{inductive pattern宣言}（\emph{inductive pattern declaration}）\index{inductivepatternせんげん@inductive pattern宣言}で宣言しておく必要がある．

inductive pattern宣言は，パターンコンストラクタの型を，データコンストラクタやマッチャーの定義から独立して宣言するための構文である．
たとえば，リスト型\lstinline{[a]}に対するパターンコンストラクタは以下のように宣言される．
\begin{lstlisting}[language=egison]
inductive pattern [a] :=
  | []
  | (::) a [a]
  | (++) [a] [a]
\end{lstlisting}
この宣言は，パターンコンストラクタ\lstinline{[]}，\lstinline{::}，\lstinline{++}がいずれもリスト型\lstinline{[a]}を対象とするパターンコンストラクタであることを宣言している．
各パターンコンストラクタの後に並ぶ型は，そのパターンコンストラクタの引数の型，すなわちパターンマッチの成功時にネクスト・ターゲットとして渡される値の型を表す．

inductive pattern宣言の重要な特徴は，\lstinline{::}というパターンコンストラクタが\lstinline{list}マッチャー，\lstinline{multiset}マッチャー，\lstinline{set}マッチャーなど，複数のマッチャーにわたって共有される点にある．それぞれのマッチャーは\lstinline{::}に対して異なる分解方法を実装するが，そのパターンコンストラクタの型はinductive pattern宣言によって一元的に管理される．
\begin{lstlisting}[language=egison,numbers=none]
matchAll [1, 2, 3] as list integer with
  | $x :: $xs -> (x, xs)
-- [(1,[2,3])]

matchAll [1, 2, 3] as multiset integer with
  | $x :: $xs -> (x, xs)
-- [(1,[2,3]),(2,[1,3]),(3,[1,2])]
\end{lstlisting}
同じパターン\lstinline{$x :: $xs}でも，使うマッチャーによって意味（分解方法）が変わる．

このinductive pattern宣言とマッチャー定義の関係は，型クラスとインスタンス定義の関係に類似している．
型クラス宣言がメソッドの型を宣言し，インスタンス定義がその具体的な実装を提供するのと同様に，inductive pattern宣言がパターンコンストラクタの型を宣言し，マッチャー定義がそのパターンコンストラクタに対する具体的な分解方法を定義する．

\medskip
\begin{center}
\begin{tabular}{ll}
\hline
\textbf{型クラス} & \textbf{パターンマッチング} \\
\hline
型クラス宣言 & inductive pattern宣言 \\
インスタンス定義 & マッチャー定義 \\
メソッド名 & パターンコンストラクタ名 \\
メソッドの型 & パターンコンストラクタの型 \\
インスタンスディスパッチ（静的） & マッチャーディスパッチ（動的） \\
\hline
\end{tabular}
\end{center}
\medskip

ただし，型クラスのインスタンス選択がコンパイル時に静的に決定されるのとは異なり，マッチャーの選択は\lstinline{match}式や\lstinline{matchAll}式においてユーザーが実行時に明示的に指定する（動的選択）．

\section{マッチャーの定義の基礎}\label{matcher-basic}

本節は，非自由データ型のなかでもっとも単純なunordered pair（要素の順番を無視するペア）のマッチャー定義をみていくことにより，マッチャー定義の方法の概略を解説する．

まずは，整数のunordred pairに対するマッチャー\lstinline{unorderedIntegerPair}マッチャーの挙動を確認しよう．
\lstinline{unorderedIntegerPair}に対して，\lstinline{pair}パターンを使うと以下のように2通りの要素の順番を両方パターンマッチする．
\begin{lstlisting}[language=egison,numbers=none]
matchAll (1, 2) as unorderedIntegerPair with
| pair $x $y -> (x, y)
-- [(1,2),(2,1)]
\end{lstlisting}
そのおかげで，\lstinline{pair}の第一引数がターゲットの2番目の要素である場合にも，パターンマッチに成功する．
\begin{lstlisting}[language=egison,numbers=none]
matchAll (1, 2) as unorderedIntegerPair with
| pair #2 $y -> y
-- [1]
\end{lstlisting}
パターンがパターン変数であった場合には，ターゲットをそのまま束縛する．
\begin{lstlisting}[language=egison,numbers=none]
matchAll (1, 2) as unorderedIntegerPair with
| $p -> p
-- [(1,2)]
\end{lstlisting}

上記のような動作をする\lstinline{unorderedIntegerPair}マッチャーは以下のように定義できる．
まず，\lstinline{pair}パターンコンストラクタをinductive pattern宣言で宣言しておく必要がある．
\begin{lstlisting}[language=egison]
inductive pattern (Integer, Integer) :=
  | pair Integer Integer
\end{lstlisting}
この宣言は，\lstinline{pair}が整数のペア型\lstinline{(Integer, Integer)}を対象とするパターンコンストラクタで，\lstinline{Integer}型の引数を2つ取ることを宣言している．
次に，型注釈により整数のペアに対するマッチャーであることを示しながら，\lstinline{unorderedIntegerPair}マッチャーを以下のように定義する．
\begin{lstlisting}[language=egison]
def unorderedIntegerPair : Matcher (Integer, Integer) := matcher
  | pair $ $ as (integer, integer) with
    | ($x, $y) -> [(x, y), (y, x)]
  | $ as something with
    | $tgt -> [tgt]
\end{lstlisting}
\lstinline{matcher}は，マッチャーを生成するための組み込み構文である．
\begin{lstlisting}[language=egison,numbers=none]
matcher
| `原始パターンパターン` as `ネクスト・マッチャー式` with
  | `原始データパターン` -> `ボディ`
  ...
...
\end{lstlisting}
\lstinline{matcher}式は，複数のマッチャー節（\emph{matcher clause}）をとる．
マッチャー節は，\emph{原始パターンパターン}（\emph{primitive pattern-pattern}）\index{げんしぱたーんぱたーん@原始パターンパターン}，ネクスト・マッチャー式（\emph{next matcher expression}），複数の原始マッチ節（\emph{primitive match clause}）からなる．
原始マッチ節は，\emph{原始データパターン}（\emph{primitive data pattern}）\index{げんしでーたぱたーん@原始データパターン}とボディからなる．

\lstinline{unorderedIntegerPair}の1つ目のマッチャー節（2-3行目）をみていく．
まず，このマッチャー節の原始パターンパターンは\lstinline{pair $ $}である．
原始パターンパターンはパターンに対するパターンである．
この原始パターンパターンは\lstinline{pair}パターンにマッチし，このマッチャー節は\lstinline{pair}パターンに対するパターンマッチの方法を定義している．
\lstinline{pair}のように，原始パターンパターン中にあらわれるパターンコンストラクタは，原始パターンパターンコンストラクタ（\emph{primitive pattern-pattern constructor}）と呼ばれる．
\lstinline{pair}パターンの引数としてあらわれている\lstinline{$}はパターン・ホール（\emph{pattern hole}）と呼ばれる原始パターンパターンの構成要素で，任意のパターンがマッチする．
パターン・ホールにマッチしたパターンは，ネクスト・パターン（\emph{next pattern}）と呼ばれる．
原始パターンパターンのパターンマッチは，一般の関数型プログラミング言語の代数的データ型に対するパターンマッチとほぼ同じようにおこなわれる．
次に，このマッチャー節のネクスト・マッチャー式は，\lstinline{(integer, integer)}である．
これは，上記のパターン・ホールに束縛された2つのパターン（\lstinline{pair}パターンの引数）がそれぞれ\lstinline{integer}マッチャーを使って再帰的にパターンマッチされることを表現している．
このマッチャー節は，1つの原始マッチ節（3行目）をとる．
原始データパターンは，ターゲットとパターンマッチされる．
原始データパターンのパターンマッチも，一般の関数型プログラミング言語の代数的データ型に対するパターンマッチとほぼ同じようにおこなわれる．
原始マッチ節のボディは，ターゲットの分解結果を返す．
\lstinline{(x, y)}と\lstinline{(y, x)}はネクスト・ターゲット（\emph{next target}）と呼ばれ，ネクスト・パターンとネクスト・マッチャーを使って再帰的にパターンマッチされる．

2つ目のマッチャー節（4-5行目）の原始パターンパターンは，パターン・ホールである．
パターン・ホールは任意のパターンにマッチするため，1つ目のマッチャー節でマッチしなかったパターンはすべて2つ目のマッチ節で処理される．
1つ目のマッチャー節でマッチしない\lstinline{unorderedIntegerPair}に対するパターンは，ワイルドカードかパターン変数のみであるので，このマッチャー節はこれらのパターンを処理する．
このマッチャー節は，パターンとターゲットを変えず，ただマッチャーを\lstinline{something}に変換するだけである．
パターンとターゲットは\lstinline{something}を使ってパターンマッチされる．
\lstinline{something}は唯一の組み込みマッチャーであり，以下の3種類のパターンを処理する．
\begin{itemize}
  \item \textbf{ワイルドカード}：つねにパターンマッチに成功し，束縛は生じない．
  \item \textbf{パターン変数} \lstinline{$x}：パターンマッチに成功し，ターゲットの値を\lstinline{x}に束縛する．
  \item \textbf{値パターン} \lstinline{#e}：式\lstinline{e}を評価した値とターゲットの値を\emph{組み込み構造的同値判定}（データコンストラクタの構造による等価比較）で比較し，等しければパターンマッチに成功する．
\end{itemize}
値パターンを\lstinline{something}が直接処理するようになったのは，静的型付けの導入にともなう設計変更である（詳しくは後述の\ref{matcher-eq}節を参照）．

以下の\lstinline{unorderedPair}のように，ペアの要素のデータ型に対するマッチャーを受け取るunordered pairに対するマッチャーを定義することができる．
\begin{lstlisting}[language=egison,numbers=none]
matchAll (1, 2) as unorderedPair integer with
| pair $x $y -> (x, y)
-- [(1,2),(2,1)]
\end{lstlisting}
そのためには，まずinductive pattern宣言を型パラメータを使った汎用的な形に変更する．
\begin{lstlisting}[language=egison,numbers=none]
inductive pattern {a} (a, a) :=
  | pair a a
\end{lstlisting}
そして，\lstinline{unorderedPair}をマッチャーを引数にとってマッチャーを返す関数として定義すればよい．
\lstinline{pair}パターンに対するマッチャー節のネクスト・マッチャー式で\lstinline{unorderedPair}の引数である\lstinline{m}を指定している．
型注釈により、型パラメータ\lstinline|{a}|と引数のマッチャー\lstinline|(m: Matcher a)|を受け取り、ペア型\lstinline|(a, a)|に対するマッチャーを返すことを示している．
\begin{lstlisting}[language=egison,numbers=none]
def unorderedPair {a} (m: Matcher a) : Matcher (a, a) := matcher
  | pair $ $ as (m, m) with
    | ($x, $y) -> [(x, y), (y, x)]
  | $ as something with
    | $tgt -> [tgt]
\end{lstlisting}

\section{\texttt{eq}マッチャーの定義}\label{matcher-eq}

\lstinline{eq}マッチャーも\lstinline{matcher}式でユーザーが定義することができる．
\lstinline{eq}マッチャーは値パターンとワイルドカード・パターン変数のみを処理するため，特別なパターンコンストラクタを持たない．
したがって，\lstinline{eq}マッチャー自体にはinductive pattern宣言は不要である．
型注釈により，型制約\lstinline|{Eq a}|を持つ型パラメータ\lstinline{a}に対するマッチャーであることを示している．
\begin{lstlisting}[language=egison]
def eq {Eq a} : Matcher a := matcher
  | #$val as () with
    | $tgt -> if val == tgt then [()] else []
  | $ as something with
    | $tgt -> [tgt]
\end{lstlisting}
\lstinline{eq}マッチャーの1つ目のマッチャー節（2-3行目）は値パターンに対するマッチャー節である．
\lstinline{#$val}は値パターンにマッチする原始パターンパターンである．
このような原始パターンパターンは原始値パターン（\emph{primitive value pattern}）と呼ばれる．
変数\lstinline{val}には値パターンの中身が束縛される．
このマッチャー節の原始パターンパターンは，パターン・ホールを含んでいないため，ネクスト・マッチャー式は空タプルである．
原始マッチ節で，値パターンの中身とターゲットの値が等しいかどうかチェックしている．

\paragraph{\texttt{something}と\texttt{eq}の違い．}
\lstinline{something}は値パターンを組み込みの構造的同値判定で処理するのに対し，\lstinline{eq}は\lstinline{Eq}型クラスのメソッドであるユーザ定義の\lstinline{==}演算子を使って比較する．
この2つの使い分けは，マッチャーを引数にとるマッチャー（マッチャーパラメータ化マッチャー）の設計と深く関わっている．

たとえば，\lstinline{multiset}のような引数マッチャー\lstinline{m}を受け取るマッチャーで値パターンを処理するとき，再帰的な比較の基底ケースをどのマッチャーに委ねるかが問題になる．
もし\lstinline{something}が値パターンを処理できないとすると，\lstinline{eq}マッチャーに委ねるしかなく，\lstinline{eq}は\lstinline{Eq}型クラス制約を必要とするため，\lstinline{multiset}の型が
\lstinline|multiset : {Eq a} => Matcher a -> Matcher [a]|
のように\lstinline{Eq}制約を必要とするものになってしまう．
しかし\lstinline{multiset}の多くの用途は値パターンを使わないため，この制約は過度に制限的である．

この問題を解決するために，\lstinline{something}に組み込みの構造的同値判定による値パターン処理を組み込んだ．
これにより，マッチャーパラメータ化マッチャーは，型クラス制約を課すことなく\lstinline{something}を基底ケースとして使え，
\lstinline|multiset : Matcher a -> Matcher [a]|
という制約のない型を持つことができる．
一方，ユーザ定義の等値比較（\lstinline{==}）を使いたい場合には，明示的に\lstinline{eq}マッチャーを指定することで対応できる．

原始パターンパターンは，ここまでにでてきた3つの構成要素であるパターン・ホール，原始パターンパターンコンストラクタの適用，原始値パターンと，原始ワイルドカードからなる．
\begin{lstlisting}[language=egison,numbers=none]
`原始パターンパターン` ::= $ | c `原始パターンパターン`* | #$`ID` | _
\end{lstlisting}
原始ワイルドカードについては，\ref{matcher-wildcard}節で用例を紹介する．

原始データパターンは，ワイルドカード，パターン変数，原始データパターンコンストラクタの適用からなる．
\begin{lstlisting}[language=egison,numbers=none]
`原始データパターン` ::= _ | $`ID` | C `原始データパターン`*
\end{lstlisting}

\section{\texttt{multiset}マッチャーの定義}\label{matcher-multiset}

本節は，\lstinline{multiset}マッチャーの定義を解説する．
\lstinline{multiset}マッチャーは，\ref{matcher-inductive}節で示したリスト型\lstinline{[a]}のinductive pattern宣言で宣言されたパターンコンストラクタ（\lstinline{[]}，\lstinline{::}，\lstinline{++}）を処理するマッチャーである．
型注釈により、型パラメータ\lstinline|{a}|と引数のマッチャー\lstinline|(m: Matcher a)|を受け取り、リスト型\lstinline|[a]|に対するマッチャーを返すことを示している．
\begin{lstlisting}[language=egison]
def multiset {a} (m: Matcher a) : Matcher [a] := matcher
  | [] as () with
    | $tgt -> match tgt as (mutiset m) with
                | [] -> [()]
                | _ -> []
  | $ :: $ as (m, multiset m) with
    | $tgt -> matchAll tgt as list m with
                | $hs ++ $x :: $ts -> (x, hs ++ ts)
  | #$val as () with
    | $tgt -> match (val, tgt) as (list m, multiset m) with
                | ([], []) -> [()]
                | ($x :: $xs, #x :: #xs) -> [()]
                | (_, _) -> []
  | $ as something with
    | $tgt -> [tgt]
\end{lstlisting}
1つ目と4つ目のマッチャー節はほぼ自明であるので，2つ目と3つ目のマッチャー節をみていく．

2つ目のマッチャー節（6-8行目）をみていこう．
原始プリミティブパターンは\lstinline{$ :: $}で，ネクスト・マッチャー式は\lstinline{(a, multiset a)}である．
\lstinline{a}は\lstinline{multiset}の引数のマッチャーである．
ネクスト・ターゲットは，\lstinline{matchAll}式を使って定義されている．
この\lstinline{matchAll}式は，ターゲットが\lstinline{[1,2,3]}である場合，\lstinline{[(1,[2,3]),(2,[1,3]),(3,[1,2])]}を返す．
このリストに含まれるそれぞれのネクスト・ターゲットは，ネクスト・パターンとネクスト・マッチャーを使って再帰的にパターンマッチされる．
たとえば，ネクスト・ターゲット\lstinline{1}と\lstinline{[2,3]}は，ネクスト・マッチャー\lstinline{a}と\lstinline{multiset a}を使って，コンス・パターンの第1引数と第2引数のパターンとパターンマッチされる．

3つ目のマッチ節をみていこう．
このマッチャー節の原始パターンパターンは，原始値パターン\lstinline|#$val|である．
このマッチャー節は，値パターンを処理する．
このマッチャー節は，値パターンの中身（\lstinline{val}）とターゲット（\lstinline{tgt}）が等しいかどうかくらべる．
この\lstinline{match}式の1つ目と3つ目のマッチ節は自明であるので説明を省略する．
2つ目のマッチ節のパターンは，\lstinline{val}の先頭の要素を\lstinline{tgt}から取り出す．
そして，\lstinline{val}と\lstinline{tgt}から同じ要素を1つ抜いたコレクションを\lstinline{$xs}と\lstinline{#xs}というパターンを使って再帰的にパターンマッチしている．
\lstinline{#xs}のパターンマッチに，このマッチャー節自身が再帰的に呼び出されるが，コレクションの長さが1つずつ短くなっていくため，最終的に1つ目か3つ目のマッチ節どちらかにいきつく．

\section{\texttt{soretedList}マッチャーの定義}\label{matcher-sorted}

二重にネストしたジョイン・コンス・パターンを使って$(p,p+6)$という形の素数のペアを列挙するプログラムは，後方の素数のペアになるほど列挙に時間がかかるようになる．
その理由は，二重にネストしたジョイン・コンス・パターンは，すべての素数の組み合わせを検査するためである．
たとえば，このパターンは，$(3,5)$，$(3,7)$，$(3,11)$，$(3,13)$，$(3, 17)$，$(3,19)$のような素数のペアすべてを検査する．
しかし，$(3,11)$以降の素数のペア，差が$6$より大きくなる最初の素数のペア，以降のペアについては，$(p,p+6)$であるか検査する必要はない．
\begin{lstlisting}[language=egison,numbers=none]
take 10 (matchAll primes as sortedList integer with
         | _ ++ $p :: (_ ++  #(p + 6) :: _) -> (p, p + 6))
-- [(5,11),(7,13),(11,17),(13,19),(17,23),(23,29),(31,37),(37,43),(41,47),(47,53)]
\end{lstlisting}

ソート済みリストに特化したマッチャーを使うことにより，この不必要な探索は避けることができる．
通常の\lstinline{list}マッチャーに，\lstinline|$ ++ #$px : $|を原始パターンパターンにもつマッチャー節を追加すれば，ソート済みリストに特化したマッチャーをつくることができる．
このマッチャー節が，$(p,p+6)$という形の素数のペアにマッチする二重にネストしたジョイン・コンス・パターンの計算量を$O(n^2)$から$O(n)$に減らす．
型注釈により、型制約\lstinline|{Ord a}|を持つ型パラメータ\lstinline{a}とマッチャー引数\lstinline|(m: Matcher a)|を受け取り、ソート済みリスト型\lstinline|[a]|に対するマッチャーを返すことを示している．
\begin{lstlisting}[language=egison,numbers=none]
def sortedList {Ord a} (m: Matcher a) : Matcher [a] := matcher
  | $ ++ #$px :: $ as (sortedList m, sortedList m) with
    | \$tgt -> matchAll tgt as list m with
               | loop $i (1, $n)
                   ((?(\x -> x < px) & $h_i) :: ...)
                   (#px :: $ts)
                 -> (map (\i -> h_i) [1..n], ts)
  ...
\end{lstlisting}
このようにネストしたパターンに対して直接パターンマッチのアルゴリズムを定義することによる最適化のことを，\emph{パターン・フュージョン}（\emph{pattern fusion}）\index{ぱたーんふゅーじょん@パターン・フュージョン}と呼ぶ．

\section{原始ワイルドカードによる最適化}\label{matcher-wildcard}

パターンがワイルドカードを含む場合，ワイルドカードにマッチするターゲットの計算を省くことにより，パターンマッチの処理を高速化できる．
原始パターンパターン中で，原始ワイルドカードを使うことにより，Egisonではこの最適化をおこなえる．
本節は，\ref{matcher-multiset}節で紹介した\lstinline{multiset}マッチャーについて，そのような最適化を紹介する．
このような最適化のことを，\emph{ワイルドカード最適化}（\emph{wildcard optimization}）\index{わいるどかーどさいてきか@ワイルドカード最適化}と呼ぶ．

本節で紹介する最適化の対象は，以下のようなコンス・パターンの第二引数がワイルドカードである場合である．
この場合，対象のコレクションから要素を1つ除いた残りのコレクションを計算する必要はない．
原始ワイルドカードを使うことにより，この計算を省くように\lstinline{multiset}マッチャーに記述することができる．
\begin{lstlisting}[language=egison]
matchAll [1, 2, 3, 4, 5] as multiset eq with
| $x :: _ -> x
-- [1, 2, 3, 4, 5]
\end{lstlisting}
このような最適化は，以下の2-3行目のマッチャー節を\ref{matcher-multiset}節の\lstinline{multiset}マッチャーの定義に追加すればできる．
このマッチャー節の原始パターンパターンで原始ワイルドカードが使われている．
パターン・ホールはコンス・パターンの第一引数だけであるので，ネクスト・マッチャーは\lstinline{m}だけになり，ネクスト・ターゲットは\lstinline{tgt}のそれぞれの要素であるために，\lstinline{tgt}をそのまま返すようになっている．
\begin{lstlisting}[language=egison]
def multiset {a} (m: Matcher a) : Matcher [a] := matcher
  | $ :: _ as m with
    | $tgt -> tgt
  | $ :: $ as (m, multiset m) with
    | $tgt -> matchAll tgt as list m with
              | $hs ++ $x :: $ts -> (x, hs ++ ts)
  ...
\end{lstlisting}

\ref{matcher-sorted}節の\lstinline{sortedList}マッチャーも同様の最適化ができる．
下記の2-7行目のマッチャー節は，ジョイン・パターンの第一引数を計算することを省略する役割を果たす．
\begin{lstlisting}[language=egison]
def sortedList {Ord a} (m: Matcher a) : Matcher [a] := matcher
  | _ ++ #$px :: $ as sortedList m
    | \tgt -> matchAll tgt as list m with
              | loop $i (1, _)
                  (?(\x -> x < px) :: ...)
                  (#px :: $ts)
                -> ts
  | $ ++ #$px :: $ as (sortedList m, sortedList m)
    | \tgt -> matchAll tgt as list m with
              | loop $i (1, $n)
                  ((?(\x -> x < px) & $h_i) :: ...)
                  (#px :: $ts)
                -> (map (\i -> h_i) [1..n], ts)
  ...
\end{lstlisting}

\section{\texttt{assocMultiset}マッチャー}\label{matcher-assoc}

Egisonには，多重集合に対するマッチャーとして，\lstinline{assocMultiset}が用意されている．
本節はその使い方と定義を紹介する．

多重集合は，要素のその個数の連想リストとして表現できる．
\lstinline{assocMultiset}マッチャーは，このように連想リストとして表現された多重集合に対するマッチャーである．
たとえば，以下のプログラムは$1$を$4$個，$2$を$3$個，$3$を$1$個もつ多重集合に対してパターンマッチをしている．
ターゲットの多重集合に$3$個以上現れる要素を取り出すパターンが表現されている．
\begin{lstlisting}[language=egison]
matchAll [(1, 4), (2, 3), (3, 1)] as assocMultiset eq with
  | $x ^ #3 :: $rs -> (x, rs)
-- [(1, [(1, 1), (2, 3), (3, 2)]), (2, [(1, 4), (3, 2)])]
\end{lstlisting}
